% Source: https://artofproblemsolving.com/community/c3h1442913_a_fraction_calculate_problem
% Stable identity: aops-topic-1442913-q01
Since $1+2+\cdots+k=\dfrac{k(k+1)}{2}$, the factor with last term $k$ is
\[
1-\dfrac{2}{k(k+1)}=\dfrac{k^2+k-2}{k(k+1)}=\dfrac{(k-1)(k+2)}{k(k+1)}.
\]
So, with $k$ running from $2$ to $997$,
\[
A=\prod_{k=2}^{997}\dfrac{k-1}{k}\cdot\prod_{k=2}^{997}\dfrac{k+2}{k+1}=\dfrac{1}{997}\cdot\dfrac{999}{3}=\boxed{\dfrac{333}{997}},
\]
because both products telescope: the first is $\dfrac{1}{2}\cdot\dfrac{2}{3}\cdots\dfrac{996}{997}$ and the second is $\dfrac{4}{3}\cdot\dfrac{5}{4}\cdots\dfrac{999}{998}$.
